\chapter{Conclusiones}

\label{Chapter5}

\section{Conclusiones generales}

Este trabajo abordó el diseño, desarrollo e implementación de un sistema modular de aprendizaje continuo para chatbots basados en Procesamiento de Lenguaje Natural (PLN), enfocado en ofrecer soluciones accesibles, escalables y sostenibles para pequeños emprendimientos. La arquitectura se basó en el enfoque RAG (Retrieval-Augmented Generation), permitiendo generar respuestas precisas y contextualizadas, con mejora continua sin necesidad de reentrenamiento.

\subsection{Cumplimiento de objetivos y planificación}

Los objetivos planteados fueron plenamente alcanzados: se desarrolló un pipeline funcional, de bajo costo y fácil de implementar, que incorpora mecanismos de evaluación automática y trazabilidad.

En términos de planificación, si bien se respetaron los recursos estimados (horas de trabajo y presupuesto de cómputo en la nube), los plazos debieron extenderse. El desarrollo previsto para finalizar en noviembre de 2024 concluyó en abril de 2025, y la redacción técnica finalizó en junio, debido a pausas por cuestiones personales y laborales.

\subsection{Gestión de riesgos y adaptaciones}

Durante el proceso, se presentaron dos riesgos contemplados originalmente:

\begin{itemize}
\item Retrasos por razones personales o laborales, que exigieron una replanificación flexible del cronograma, sin comprometer la calidad del resultado.
\item Limitaciones en la retroalimentación a partir de historiales, que se resolvieron exitosamente mediante la incorporación de Langfuse, lo que mejoró la trazabilidad y permitió una evaluación automática más eficaz.
\end{itemize}

\subsection{Resultados, aprendizajes y aportes}

El sistema mostró un rendimiento sólido en distintos escenarios, incluyendo consultas de múltiples pasos o fuera de contexto. Se observaron mejoras significativas en la coherencia y relevancia de las respuestas a lo largo de las distintas versiones del prototipo. Las estrategias de reformulación de preguntas, clasificación temática y recuperación semántica fueron fundamentales para lograr estos resultados sin reentrenamiento del modelo.

Desde el punto de vista formativo, este trabajo integró múltiples conceptos de la especialización en Inteligencia Artificial: diseño de arquitecturas modulares, evaluación automática, procesamiento semántico, vectores de embeddings, sistemas conversacionales y despliegue en la nube. Estas herramientas no solo permitieron construir una solución técnicamente sólida, sino también comprender en profundidad los desafíos prácticos que implica desarrollar aplicaciones basadas en PLN en entornos reales.

Los principales aportes del trabajo incluyen:

\begin{itemize}
\item Una arquitectura modular basada en DSPy y LanceDB, adaptable y de bajo mantenimiento.
\item Un sistema completo de evaluación conversacional, con métricas automáticas, análisis semántico y agrupamiento temático.
\item Un entorno desplegado en Azure Container Apps, listo para su uso en pruebas reales sin depender de infraestructura externa costosa.
\end{itemize}

\section{Próximos pasos}

Una evolución natural de este sistema consiste en integrar el módulo RAG como un agente especializado dentro de un chatbot multiagente. Por ejemplo, en un entorno gastronómico, se podrían definir distintos agentes colaborativos:

\begin{itemize}
\item Un agente de atención al cliente, encargado de responder preguntas frecuentes sobre horarios, menú o ubicación.
\item Un agente de pedidos, que registre solicitudes de clientes y las almacene en una base de datos interna.
\item Un agente CRM, conectado a sistemas externos para gestionar datos de clientes y su historial de interacciones.
\item Un agente RAG, responsable de consultas complejas o basadas en documentación, aprovechando la arquitectura desarrollada en este trabajo.
\end{itemize}

Este tipo de arquitectura distribuida y colaborativa permitiría construir asistentes conversacionales más flexibles, extensibles y adaptados a distintos dominios, con una lógica clara de mantenimiento y evolución continua.

\section{Síntesis final}

El desarrollo realizado sienta las bases para una nueva generación de chatbots inteligentes, modulares y autoevaluables. La arquitectura propuesta demuestra que es posible diseñar soluciones conversacionales robustas, sin necesidad de reentrenamiento, con bajo costo operativo y aplicabilidad directa en contextos reales. Este trabajo constituye una contribución concreta a la implementación práctica de tecnologías de PLN en entornos accesibles para pequeñas organizaciones.
